\documentclass[12pt]{article}
\usepackage[utf8]{inputenc}
\usepackage[T1]{fontenc}
\usepackage{hyperref}
\title{Kunfal Language Resources for K-12 Curriculum Development}
\author{Compiled by Claude Sonnet 5.5 (claude-sonnet-5-5) for Andr\'as Kornai}
\date{}
\begin{document}
\maketitle

\section{Introduction}
Kunfal (ISO 639-3 xuf) is a retired code. Glottolog lists it under ``Bookkeeping''. The ISO 639-3 Registration Authority merged it into Awngi (awn) in 2008, on the basis of 2005 surveys by Joswig, Mohammed and Boone. Those surveys found that Kunfal speakers understand Awngi texts without difficulty and regard it as their own language. I could not confirm the Glottocode kunf1238 in my searches. Kunfal materials are therefore best sought under Awngi, an Agaw language of Ethiopia. This dossier does not cover Awngi itself in depth.

\section{Primary Educational Resources}
I found no Kunfal-specific primers, textbooks or school programmes in the search results.

\section{Linguistic and Reference Materials}
\begin{itemize}
\item Glottolog entry for xuf, giving its retired status and the merger into Awngi. \url{https://glottolog.org/resource/languoid/iso/xuf}
\item Glottolog record for the ISO 639-3 Registration Authority 2007 change request. \url{https://glottolog.org/resource/reference/id/500029}
\item Birru, Adal and Cowley (1971), ``The Kunf\"al people and their language'', Journal of Ethiopian Studies 9(2):99--106. This is cited as the main early source. I found a citation only, not the article itself.
\end{itemize}

\section{Religious and Community Resources}
I found nothing specific to Kunfal.

\section{Digital and Audio Resources}
I found nothing specific to Kunfal.

\section{Curriculum Development Context}
Because the code has been retired, the practical route is to treat Kunfal as an Awngi variety and use Awngi educational resources. Those need a separate search under awn. No Kunfal-specific curriculum is known.

\end{document}